\section{Playfair Cipher}

% Matriks kunci Playfair "JALAN GANESHA SEPULUH".
% #1 = sel plainteks (baris/kolom), #2 = sel cipherteks (baris/kolom); baris 0 di atas.
\newcommand{\pfpanel}[2]{%
  \begin{tikzpicture}[x=0.5cm, y=0.5cm]
    \foreach \r/\c in {#1} {\fill[gray!35] (\c,-\r) rectangle ++(1,-1);}
    \foreach \r/\c in {#2} {\fill[green!35] (\c,-\r) rectangle ++(1,-1);}
    \draw (0,0) grid (5,-5);
    \foreach \r/\rowtxt in {0/{A,L,N,G,E},1/{S,H,P,U,B},2/{C,D,F,I,K},3/{M,O,Q,R,T},4/{V,W,X,Y,Z}} {
      \foreach \ch [count=\c from 0] in \rowtxt {
        \node[font=\ttfamily\small] at (\c+0.5,-\r-0.5) {\ch};
      }
    }
  \end{tikzpicture}%
}

\begin{figure}[htbp]
  \centering
  \begin{minipage}{0.3\textwidth}
    \centering
    \includegraphics[width=0.62\textwidth]{bab-04/potret-wheatstone}\\[2pt]
    \small Sir Charles Wheatstone
  \end{minipage}
  \begin{minipage}{0.3\textwidth}
    \centering
    \includegraphics[width=0.62\textwidth]{bab-04/potret-playfair}\\[2pt]
    \small Baron Lyon Playfair
  \end{minipage}
  \caption{Penemu dan promotor Playfair cipher \cite{munir_slides_itb}.}
  \label{fig:potret-playfair}
\end{figure}

Playfair cipher ditemukan oleh Sir Charles Wheatstone, namun dipromosikan oleh Baron Lyon Playfair pada tahun 1854 sehingga dinamai menurut nama Playfair (Gambar~\ref{fig:potret-playfair}) \cite{munir_slides_itb,stallings2017}.

\subsection{Konsep}

Playfair cipher termasuk kelompok \crypt{cipher poligram} (\textit{polygram cipher}): satu blok huruf plainteks disubstitusi dengan satu blok huruf cipherteks. Blok 2 huruf disebut \textit{bigram} (digram), blok 3 huruf disebut trigram, dan seterusnya. Playfair cipher mengenkripsi dalam bentuk bigram; misalnya bigram \texttt{AS} diganti dengan \texttt{RT}, dan \texttt{BY} diganti dengan \texttt{SL}.

Kunci enkripsi/dekripsi disusun sebagai matriks $5 \times 5$ yang setiap elemennya huruf alfabet. Karena hanya ada 25 elemen untuk 26 huruf, satu huruf dibuang, yaitu huruf J. Banyaknya matriks kunci yang mungkin adalah
\[
25! = 15.511.210.043.330.985.984.000.000.
\]

\begin{catatan}
\textbf{Mengapa huruf J yang dibuang?} Dalam alfabet Latin lama, I dan J tidak dibedakan; J hanyalah variasi penulisan I. Karena itu Playfair menggabungkan keduanya: setiap J pada plainteks diganti dengan I (contoh: \texttt{JALAN} $\rightarrow$ \texttt{IALAN}). Saat dekripsi, huruf I dapat berarti I atau J, dan ambiguitas kecil ini diselesaikan dari konteks kalimat. Secara teori huruf lain bisa saja dibuang, tetapi I/J paling masuk akal dan sudah menjadi konvensi.
\end{catatan}

\subsection{Pembentukan Matriks Kunci}

Kunci dapat dipilih dari kalimat yang mudah diingat, misalnya \texttt{JALAN GANESHA SEPULUH}.
\begin{enumerate}
  \item Buang huruf yang berulang dan huruf J: \texttt{ALNGESHPU}.
  \item Tambahkan huruf-huruf yang belum ada (kecuali J): \texttt{ALNGESHPUBCDFIKMOQRTVWXYZ}.
  \item Masukkan baris demi baris ke dalam bujursangkar $5 \times 5$ (Gambar~\ref{fig:playfair-kunci}).
\end{enumerate}

\begin{figure}[htbp]
  \centering
  \pfpanel{}{}
  \caption{Matriks kunci Playfair untuk kalimat kunci \texttt{JALAN GANESHA SEPULUH}.}
  \label{fig:playfair-kunci}
\end{figure}

\subsection{Praproses Plainteks}

Sebelum dienkripsi, pesan diatur sebagai berikut:
\begin{enumerate}
  \item Buang semua spasi di dalam pesan.
  \item Ganti huruf J (bila ada) dengan I.
  \item Tulis pesan dalam pasangan huruf (bigram).
  \item Jangan sampai ada bigram dengan dua huruf yang sama; jika ada, sisipkan X di tengahnya.
  \item Jika jumlah huruf ganjil, tambahkan X pada bigram terakhir.
\end{enumerate}

\begin{contoh}
Plainteks \texttt{temui ibu nanti malam} tidak mengandung huruf J. Setelah spasi dibuang dan ditulis per bigram: \texttt{te mu ii bu na nt im al am}. Bigram \texttt{ii} berhuruf sama, sehingga disisipkan X: \texttt{te mu ix ib un an ti ma la m}. Bigram terakhir hanya satu huruf, sehingga ditambah X:
\begin{center}
\texttt{te mu ix ib un an ti ma la mx}
\end{center}
\end{contoh}

\subsection{Algoritma Enkripsi dan Dekripsi}

Setiap bigram dienkripsi dengan salah satu dari tiga aturan berikut (Gambar~\ref{fig:playfair-aturan}):
\begin{enumerate}
  \item Jika kedua huruf terdapat pada \textbf{baris} yang sama, setiap huruf diganti dengan huruf di kanannya (siklik: setelah kolom terakhir kembali ke kolom pertama).
  \item Jika kedua huruf terdapat pada \textbf{kolom} yang sama, setiap huruf diganti dengan huruf di bawahnya (siklik).
  \item Jika kedua huruf tidak sebaris maupun sekolom, keduanya membentuk sudut-sudut sebuah \textbf{persegi panjang}. Huruf pertama diganti dengan huruf pada perpotongan baris huruf pertama dengan kolom huruf kedua; huruf kedua diganti dengan huruf pada sudut keempat persegi panjang tersebut.
\end{enumerate}

\begin{figure}[htbp]
  \centering
  \begin{minipage}{0.3\textwidth}
    \centering
    \pfpanel{2/1,2/3}{2/2,2/4}\\[3pt]
    \small Baris: \texttt{di} $\rightarrow$ \texttt{FK}
  \end{minipage}\hfill
  \begin{minipage}{0.3\textwidth}
    \centering
    \pfpanel{3/2,3/4}{3/3,3/0}\\[3pt]
    \small Baris (siklik): \texttt{qt} $\rightarrow$ \texttt{RM}
  \end{minipage}\hfill
  \begin{minipage}{0.3\textwidth}
    \centering
    \pfpanel{0/2,3/2}{1/2,4/2}\\[3pt]
    \small Kolom: \texttt{nq} $\rightarrow$ \texttt{PX}
  \end{minipage}

  \vspace{1em}
  \begin{minipage}{0.3\textwidth}
    \centering
    \pfpanel{3/1,4/1}{0/1}\\[3pt]
    \small Kolom (siklik): \texttt{ow} $\rightarrow$ \texttt{WL}
  \end{minipage}\hspace{0.04\textwidth}
  \begin{minipage}{0.3\textwidth}
    \centering
    \pfpanel{1/1,4/4}{1/4,4/1}\\[3pt]
    \small Persegi panjang: \texttt{hz} $\rightarrow$ \texttt{BW}
  \end{minipage}
  \caption{Tiga aturan enkripsi Playfair. Sel abu-abu adalah huruf plainteks, sel hijau adalah huruf cipherteks. Pada \texttt{ow}, huruf W sekaligus menjadi hasil enkripsi O.}
  \label{fig:playfair-aturan}
\end{figure}

\begin{contoh}
Enkripsi bigram \texttt{te mu ix ib un an ti ma la mx} dengan matriks Gambar~\ref{fig:playfair-kunci}:
\begin{center}
\small
\begin{tabular}{cll@{\hspace{2em}}cll}
  \toprule
  Bigram & Aturan & Hasil & Bigram & Aturan & Hasil \\
  \midrule
  \texttt{te} & kolom & \texttt{ZB} & \texttt{an} & baris & \texttt{LG} \\
  \texttt{mu} & persegi panjang & \texttt{RS} & \texttt{ti} & persegi panjang & \texttt{RK} \\
  \texttt{ix} & persegi panjang & \texttt{FY} & \texttt{ma} & kolom (siklik) & \texttt{VS} \\
  \texttt{ib} & persegi panjang & \texttt{KU} & \texttt{la} & baris & \texttt{NL} \\
  \texttt{un} & persegi panjang & \texttt{PG} & \texttt{mx} & persegi panjang & \texttt{QV} \\
  \bottomrule
\end{tabular}
\end{center}
Cipherteks: \texttt{ZB RS FY KU PG LG RK VS NL QV}. Sebagai ilustrasi, \texttt{t} dan \texttt{e} berada pada kolom terakhir sehingga masing-masing diganti huruf di bawahnya (\texttt{Z} dan \texttt{B}); \texttt{m} dan \texttt{u} membentuk persegi panjang sehingga \texttt{m} diganti \texttt{R} (baris M, kolom U) dan \texttt{u} diganti \texttt{S} (baris U, kolom M).
\end{contoh}

Algoritma \textbf{dekripsi} adalah kebalikan dari algoritma enkripsi:
\begin{enumerate}
  \item Jika kedua huruf terdapat pada baris yang sama, setiap huruf diganti dengan huruf di kirinya.
  \item Jika kedua huruf terdapat pada kolom yang sama, setiap huruf diganti dengan huruf di atasnya.
  \item Jika keduanya membentuk persegi panjang, aturannya sama seperti enkripsi.
  \item Buang huruf X yang tidak bermakna.
\end{enumerate}

\begin{catatan}
\textbf{Kakas daring.} Playfair cipher dapat dicoba di \url{https://planetcalc.com/7751/}, yang juga menampilkan matriks kuncinya.
\end{catatan}

\subsection{Implementasi dalam C}

Kode~\ref{lst:playfair-c} memisahkan tiga tahap: membangun matriks kunci (huruf J dibuang), praproses plainteks, dan enkripsi per bigram.

\begin{ccode}[caption={Playfair cipher dalam C}, label={lst:playfair-c}]
#include <stdio.h>
#include <ctype.h>

void build_key_square(const char *key, char sq[5][5]) {
    const char *alpha = "ABCDEFGHIJKLMNOPQRSTUVWXYZ";
    int used[26] = {0}, n = 0;
    used['J' - 'A'] = 1;                       /* huruf J dibuang */
    for (int pass = 0; pass < 2; pass++) {
        const char *s = (pass == 0) ? key : alpha;
        for (int i = 0; s[i] != '\0'; i++) {
            char ch = toupper((unsigned char)s[i]);
            if (!isalpha((unsigned char)ch) || used[ch - 'A']) continue;
            used[ch - 'A'] = 1;
            sq[n / 5][n % 5] = ch;
            n++;
        }
    }
}

/* buang non-huruf, J -> I, sisipkan X di antara huruf kembar, genapkan */
void playfair_prepare(const char *msg, char *out) {
    int n = 0;
    for (int i = 0; msg[i] != '\0'; i++) {
        char ch = toupper((unsigned char)msg[i]);
        if (!isalpha((unsigned char)ch)) continue;
        if (ch == 'J') ch = 'I';
        if (n % 2 == 1 && out[n - 1] == ch) out[n++] = 'X';
        out[n++] = ch;
    }
    if (n % 2 == 1) out[n++] = 'X';
    out[n] = '\0';
}

static void find_pos(char sq[5][5], char ch, int *r, int *c) {
    for (int i = 0; i < 25; i++)
        if (sq[i / 5][i % 5] == ch) { *r = i / 5; *c = i % 5; return; }
}

void playfair_encrypt(const char *pt, char sq[5][5], char *ct) {
    int i, r1, c1, r2, c2;
    for (i = 0; pt[i] != '\0'; i += 2) {
        find_pos(sq, pt[i], &r1, &c1);
        find_pos(sq, pt[i + 1], &r2, &c2);
        if (r1 == r2) {                        /* sebaris: geser kanan */
            c1 = (c1 + 1) % 5;  c2 = (c2 + 1) % 5;
        } else if (c1 == c2) {                 /* sekolom: geser bawah */
            r1 = (r1 + 1) % 5;  r2 = (r2 + 1) % 5;
        } else {                               /* persegi panjang: tukar kolom */
            int t = c1;  c1 = c2;  c2 = t;
        }
        ct[i] = sq[r1][c1];
        ct[i + 1] = sq[r2][c2];
    }
    ct[i] = '\0';
}

int main(void) {
    char sq[5][5], pt[64], ct[64];
    build_key_square("JALAN GANESHA SEPULUH", sq);
    playfair_prepare("temui ibu nanti malam", pt);
    playfair_encrypt(pt, sq, ct);
    printf("%s -> %s\n", pt, ct);   /* TEMUIXIBUNANTIMALAMX -> ZBRSFYKUPGLGRKVSNLQV */
    return 0;
}
\end{ccode}

\subsection{Pengayaan: Kriptanalisis Playfair Cipher}

Karena ada 26 huruf, terdapat $26 \times 26 = 676$ bigram, sehingga identifikasi bigram individual lebih sukar daripada huruf tunggal. Playfair cipher memang sulit dipecahkan dengan analisis frekuensi huruf tunggal, tetapi ukuran poligramnya hanya dua huruf, sehingga cipher ini tetap tidak aman: ia dapat dipecahkan dengan \textbf{analisis frekuensi bigram} bila cipherteks cukup panjang. Dalam bahasa Inggris, bigram TH dan HE paling sering muncul.

Kelemahan lain: sebuah bigram dan kebalikannya (misalnya AB dan BA) didekripsi menjadi pola plainteks yang juga terbalik (misalnya RE dan ER). Bahasa Inggris memiliki banyak kata dengan bigram terbalik seperti \textit{REceivER} dan \textit{DEpartED}, sehingga pola ini membantu kriptanalis.

\begin{contoh}
Misalkan bigram tersering di cipherteks adalah BM, sedangkan bigram tersering bahasa Inggris adalah TH, sehingga diduga TH $\rightarrow$ BM. Jika TH dienkripsi dengan aturan persegi panjang, maka T sebaris dengan B dan H sebaris dengan M. Jika dengan aturan kolom, B berada di bawah T dan M di bawah H; jika dengan aturan baris, B di kanan T dan M di kanan H. Dari hipotesis-hipotesis seperti ini, kriptanalis mulai membangun matriks parsial $5 \times 5$ dan mengujinya pada bigram lain.
\end{contoh}

